% R32 section 6: the X gate, compared with H (input by R32_floquet_spectra.tex)
% Numbers: tab_r32_X_N14.tex (spectral.xgate.write_x_table), results_spectra/summary_X*.json,
% results_spectra/X_cycle_growth.json (spectral.xgate.cycle_growth).
With $V=X$ every rule is a classical reversible cellular automaton. $U$ permutes the basis, so its spectrum is fixed
by the cycle structure alone: a cycle of length $L$ contributes the $L$-th roots of unity, once each. The eigenvectors
are the Fourier modes of the cycles,
\[
  |q,c\rangle = L^{-1/2}\sum_{t=0}^{L-1} e^{2\pi i q t/L}\,U^t|x_c\rangle .
\]
Level statistics therefore measure orbit arithmetic, not chaos, and are not reported. The informative quantities are
the cycle-length distribution and the entanglement of the Fourier modes. A mode is supported on $L$ basis states, so
its Schmidt rank is at most $L$ and $S_\mathrm{vN}\le\ln L$: short cycles force low entanglement.

\paragraph{Data and checks.} All sixteen rules were run with the same Julia pipeline, on the open chain for $N=6\ldots14$
and on rings for even $N=6\ldots14$: 224 runs and 40~minutes in total, with no failure. For every run, the stored
quasienergies equal the union of the roots of unity over the cycles to $2\times10^{-14}$. The cycles were computed
independently in Python from the permutation (\fn{spectral.xgate.cycle\_lengths}). A test checks this against the dense
permutation matrix at $N=8$, for both boundaries. The cycle growth in Fig.~\ref{fig:xcycles} was computed in Python up
to $N=18$.

\subsection{What the sixteen automata are, classically}
Writing $m,n$ for the left and right neighbour, the X-gate update is $x_i\to x_i\oplus f(m,n)$, with $f$ the firing
pattern:
\begin{itemize}
\item \textbf{Linear or affine} ($f$ of degree one): the XOR rules 150 ($f=m\oplus n$) and 105 ($1\oplus m\oplus n$),
  and all one-sided rules (60: $f=m$; 102: $f=n$; 153: $1\oplus n$; 195: $1\oplus m$). These are solvable by linear
  algebra over $\mathbb F_2$. X-102 (flip iff the right neighbour is 1) is exactly the deterministic Floquet quantum East
  model, whose quench dynamics are solved exactly~\cite{bertini2024det}; 60 is its mirror image.
\item \textbf{Known interacting integrable automata.} 54 fires iff at least one neighbour is 1: this is RCA54, the
  interacting integrable Rule~54 with solitons~\cite{buca2021}. 201 fires iff both neighbours are 0: this is the
  Floquet-PXP automaton~\cite{wilkinson2020}. On the ring the global spin flip is a symmetry of the X automata. It maps
  147 onto 54, 108 onto 201, 156 onto 198 and 57 onto 99. With $H$ it is not a symmetry, and the obc0 boundary (ghost sites
  fixed to 0) breaks it too: for example, 108 and 201 have 4868 and 3492 cycles on the open chain at $N=14$, but 4156
  each on the ring.
\item \textbf{Trivial}: 204 ($U=1$) and 51 ($U=X^{\otimes N}$).
\item \textbf{Nonlinear and not known to be integrable}: only 57 and 99 (fire unless $(m,n)=(0,1)$, or $(1,0)$).
\end{itemize}

\subsection{Cycle structure: only 57/99 have exponentially long orbits}
\begin{figure}[h]
\centering
\ifgraphic{../../figures/r32/fig_r32_X_cycles.pdf}{\textwidth}
\caption{X gate: (a) the longest cycle and (b) the mean cycle length seen by a uniformly random basis state,
$\langle L\rangle_\mathrm{state}=\sum_c L_c^2/2^N$, over $2^N$. Filled: open chain; open symbols: ring. Computed from the
permutation up to $N=18$.}
\label{fig:xcycles}
\end{figure}
On the open chain, the longest cycle of 57/99 holds a \emph{finite fraction} of the Hilbert space: $129\,880$ of $2^{18}$
states at $N=18$. A random state lies on a cycle of average length $0.22$ to $0.50$ times $2^N$ for every $N$ from 6 to 18.
Classically, this is an ergodic orbit structure. Every other rule has orbits that grow at most linearly:
\begin{itemize}
\item RCA54: longest cycle $3N-5$ for 54 (147, which is RCA54 only on the ring: 26 at $N=14$);
\item XOR: $N+1$;
\item one-sided: 16 or 32, independent of $N$ (a period that is a power of two, as for linear automata);
\item Floquet-PXP (201, 108): slow growth, 187 states at $N=18$.
\end{itemize}
On the ring, the cycles of 57/99 still grow exponentially, but more slowly (4428 states at $N=18$), and their share of the
space falls to $\langle L\rangle_\mathrm{state}/2^N=0.007$ at $N=14$. The classically ergodic orbit of 57/99 on the open
chain needs the boundary. The reason is a conservation law: on the ring X-57 conserves the number of domain walls, which the
fixed-$0$ ends of the open chain create and destroy (Sec.~\ref{sec:answers}). This is the mirror image of the one-sided rules with $H$, which need the ring to become chaotic.

\subsection{Eigenstate entanglement: H against X at one N}
\begin{table}[h]
\centering\scriptsize
\begin{tabular}{llrrrrr|rrrrr}
\toprule
 & & \multicolumn{5}{c|}{open chain} & \multicolumn{5}{c}{ring} \\
$k$ & $W$ & cycles & longest & $\langle L\rangle/2^N$ & $\bar S_X$ & $\bar S_H$ & cycles & longest & $\langle L\rangle/2^N$ & $\bar S_X$ & $\bar S_H$ \\
\midrule
0 & 204 & 16384 & 1 & 6.1e-05 & 0.00 & 0.00 & 16384 & 1 & 6.1e-05 & 0.00 & 0.00 \\
1 & 108 & 4868 & 33 & 4.3e-04 & 0.11 & 0.06 & 4156 & 77 & 5.3e-04 & 0.17 & 0.12 \\
1 & 156 & 2280 & 60 & 7.9e-04 & 0.14 & 0.02 & 2202 & 35 & 6.5e-04 & 0.19 & 0.04 \\
1 & 198 & 2280 & 60 & 7.9e-04 & 0.14 & 0.02 & 2202 & 35 & 6.5e-04 & 0.19 & 0.04 \\
1 & 201 & 3492 & 55 & 6.9e-04 & 0.14 & 0.07 & 4156 & 77 & 5.3e-04 & 0.17 & 0.11 \\
2 one-sided & 60 & 1044 & 16 & 9.7e-04 & 0.44 & 0.47 & 300 & 63 & 3.8e-03 & 0.84 & 0.89 \\
2 one-sided & 102 & 1044 & 16 & 9.7e-04 & 0.44 & 0.47 & 300 & 63 & 3.8e-03 & 0.84 & 0.89 \\
2 one-sided & 153 & 1024 & 16 & 9.8e-04 & 0.47 & 0.66 & 300 & 63 & 3.8e-03 & 0.84 & 0.92 \\
2 one-sided & 195 & 1024 & 16 & 9.8e-04 & 0.47 & 0.66 & 300 & 63 & 3.8e-03 & 0.84 & 0.92 \\
2 XOR & 105 & 1096 & 15 & 9.1e-04 & 0.59 & 0.52 & 1172 & 14 & 8.5e-04 & 0.58 & 0.75 \\
2 XOR & 150 & 1096 & 15 & 9.1e-04 & 0.59 & 0.51 & 2344 & 7 & 4.3e-04 & 0.43 & 0.57 \\
3 & 54 & 632 & 37 & 1.6e-03 & 0.71 & 0.89 & 732 & 91 & 2.1e-03 & 0.65 & 0.82 \\
3 & 57 & 34 & 4288 & 2.4e-01 & 0.98 & 0.98 & 500 & 462 & 7.0e-03 & 0.69 & 0.99 \\
3 & 99 & 34 & 4288 & 2.4e-01 & 0.98 & 0.98 & 500 & 462 & 7.0e-03 & 0.69 & 0.99 \\
3 & 147 & 632 & 26 & 1.6e-03 & 0.70 & 0.93 & 732 & 91 & 2.1e-03 & 0.65 & 0.93 \\
4 & 51 & 8192 & 2 & 1.2e-04 & 0.16 & 0.00 & 8192 & 2 & 1.2e-04 & 0.16 & 0.00 \\
\bottomrule
\end{tabular}

\caption{X gate at $N=14$, both boundaries: number of cycles, the longest, $\langle L\rangle_\mathrm{state}/2^N$, and the
mean half-chain entanglement of all eigenstates over the Page value ($\bar S_X$). $\bar S_H$ is the same mean for the H
gate, over its non-degenerate eigenstates (degenerate ones have basis-dependent entanglement).}
\label{tab:x}
\end{table}

\begin{figure}[h]
\centering
\ifgraphic{../../figures/r32/fig_r32_HX_compare_N14.pdf}{\textwidth}
\caption{H against X at $N=14$, open chain. (a) Mean eigenstate entanglement over the Page value. (b) X gate:
$\langle L\rangle_\mathrm{state}/2^N$.}
\label{fig:hx}
\end{figure}

Table~\ref{tab:x} and Fig.~\ref{fig:hx} give the following picture:
\begin{itemize}
\item \textbf{57/99:} on the open chain, X and H are \emph{equally} thermal ($0.98$ of Page). The cycle Fourier modes of an
  ergodic orbit look like random states. On the ring, the X value drops to $0.69$ with the shorter cycles, while H stays
  at $0.99$.
\item \textbf{54/147 (RCA54):} X gives $0.70$, bounded by its short orbits ($\ln 37=3.6$ against Page 4.35). H gives
  $0.89$ and $0.93$, and its spectrum is GOE (Sec.~\ref{sec:results}). \textbf{The Hadamard turns an interacting integrable
  classical automaton into a chaotic quantum circuit.}
\item \textbf{One-sided:} $0.44$ to $0.47$ (X) against $0.47$ to $0.66$ (H) on the open chain. On the ring both rise
  (X: $0.84$; H: $0.89$ to $0.92$), but only H is chaotic there; X stays affine, with period at most a few hundred.
\item \textbf{XOR:} X ($0.59$) is slightly \emph{above} H ($0.52$) on the open chain; both are arithmetic.
\item \textbf{$k=1$:} both are low (X $0.11$ to $0.14$; H $0.02$ to $0.07$). These are fragmented rooms, whatever the gate.
\end{itemize}

\subsection{What the comparison says}
\begin{enumerate}
\item Quantum chaos with $H$ does \emph{not} require classical chaos with $X$. Two of the four chaotic $k=3$ rules, 54 and
  147, are the interacting integrable RCA54 with $X$, yet GOE-chaotic with $H$. On the ring, 57/99 keeps its CUE chaos
  with $H$, although its X orbits no longer fill a finite share of the space. What the chaotic rules share is the firing
  count $k=3$, not the classical dynamics. Deforming classical automata by non-commuting gates is the setting of
  Ref.~\cite{iadecola2020}; for these rules, the deformation by $H$ yields chaos.
\item Conversely, the only classically ergodic automaton (57/99 on the open chain) is also chaotic with $H$.
\item With $X$, every rule splits into cycles, so ``fragmentation'' is total and carries no information. With $H$, the
  fragmentation is a property of the rule class: exponential at $k=1$ on both boundaries, and boundary-driven for the
  one-sided rules.
\end{enumerate}
